\documentclass[11pt,a4paper]{article}
\usepackage[utf8]{inputenc}
\usepackage{amsmath,amsfonts,amssymb}
\usepackage{graphicx}
\usepackage{booktabs}
\usepackage{longtable}
\usepackage{siunitx}
\usepackage{xcolor}
\usepackage{hyperref}
\usepackage{caption}
\usepackage{listings}
\usepackage{xcolor}

\lstset{
    language=Python,
    basicstyle=\ttfamily\small,
    keywordstyle=\color{blue},
    commentstyle=\color{gray},
    stringstyle=\color{orange},
    numbers=left,
    numberstyle=\tiny,
    frame=single,
    breaklines=true
}

\title{DINOv3 Deepfake Detection: Supplementary Materials}
\author{}
\date{}

\begin{document}
\maketitle

\section{Detailed Architecture Specifications}

\subsection{Model Configuration}

\begin{table}[h]
\centering
\caption{DINOv3 Encoder Configuration}
\begin{tabular}{ll}
\toprule
\textbf{Parameter} & \textbf{Value} \\
\midrule
Model Name & vit\_base\_patch16\_dinov3.lvd1689m \\
Architecture & Vision Transformer (ViT) \\
Patch Size & 16 $\times$ 16 \\
Embedding Dimension & 768 \\
Number of Layers & 12 \\
Number of Heads & 12 \\
Input Resolution & 224 $\times$ 224 \\
Parameters & 86M (frozen: $\sim$85.5M, trainable: $\sim$0.5M) \\
\bottomrule
\end{tabular}
\end{table}

\subsection{Training Hyperparameters}

\begin{table}[h]
\centering
\caption{Training Configuration}
\begin{tabular}{ll}
\toprule
\textbf{Hyperparameter} & \textbf{Value} \\
\midrule
Optimizer & AdamW \\
Learning Rate & $10^{-4}$ \\
Weight Decay & $10^{-4}$ \\
Batch Size & 16 \\
Epochs & 20 \\
Scheduler & Cosine Annealing (T\_max=20) \\
Loss Weights & $\lambda_{\text{CE}}=1.0$, $\lambda_{\text{metric}}=0.5$ \\
Angular Margin & $m=0.6$ \\
Gradient Scaling & Automatic Mixed Precision (AMP) \\
\bottomrule
\end{tabular}
\end{table}

\section{Dataset Statistics}

\begin{table}[h]
\centering
\caption{FaceForensics++ (FF++) Dataset Composition}
\begin{tabular}{lccc}
\toprule
\textbf{Manipulation Method} & \textbf{Train Videos} & \textbf{Val Videos} & \textbf{Total Frames} \\
\midrule
Original (Real) & 720 & 140 & $\sim$1.4M \\
Deepfakes & 720 & 140 & $\sim$1.4M \\
Face2Face & 720 & 140 & $\sim$1.4M \\
FaceSwap & 720 & 140 & $\sim$1.4M \\
NeuralTextures & 720 & 140 & $\sim$1.4M \\
\midrule
\textbf{Total} & 3,600 & 700 & $\sim$7M \\
\bottomrule
\end{tabular}
\end{table}

\begin{table}[h]
\centering
\caption{Celeb-DF (Validation) Dataset}
\begin{tabular}{lc}
\toprule
\textbf{Category} & \textbf{Count} \\
\midrule
Real Videos & 590 \\
Fake Videos & 5,639 \\
Total Duration & $\sim$2.5 hours \\
Unique Identities & 59 \\
\bottomrule
\end{tabular}
\end{table}

\section{Per-Method Performance Breakdown}

\begin{table}[h]
\centering
\caption{FF++ Test Set Performance by Manipulation Type}
\label{tab:per-method}
\begin{tabular}{lccc}
\toprule
\textbf{Manipulation} & \textbf{Accuracy} & \textbf{Precision} & \textbf{Recall} \\
\midrule
Deepfakes & 0.94 & 0.93 & 0.95 \\
Face2Face & 0.89 & 0.88 & 0.90 \\
FaceSwap & 0.91 & 0.90 & 0.92 \\
NeuralTextures & 0.87 & 0.86 & 0.88 \\
\midrule
\textbf{Average} & \textbf{0.90} & \textbf{0.89} & \textbf{0.91} \\
\bottomrule
\end{tabular}
\end{table}

\section{Computational Efficiency Analysis}

\begin{table}[h]
\centering
\caption{Resource Requirements Comparison}
\begin{tabular}{lccc}
\toprule
\textbf{Method} & \textbf{GPU Memory (GB)} & \textbf{CPU Memory (GB)} & \textbf{Model Size (MB)} \\
\midrule
Xception & 2.8 & 4.2 & 88 \\
EfficientNet-B4 & 3.1 & 4.5 & 75 \\
ResNet18 & 1.2 & 2.8 & 45 \\
DINOv3 (Ours) & 2.1 & 3.5 & 330 \\
CLIP (IvyFake) & 2.5 & 4.0 & 590 \\
\bottomrule
\end{tabular}
\end{table}

\section{Failure Mode Analysis}

\begin{table}[h]
\centering
\caption{Common Failure Cases and Mitigation Strategies}
\begin{tabular}{p{4cm}p{3cm}p{6cm}}
\toprule
\textbf{Failure Case} & \textbf{Frequency} & \textbf{Mitigation} \\
\midrule
No face detected & 5-8\% & Fallback to full-frame ResNet18 detector \\
Low confidence ($<$0.95) crops & 3-5\% & Lower threshold, manual review flag \\
Heavy compression artifacts & 2-4\% & Quality-aware ensemble \\
Extreme pose/occlusion & 4-6\% & Multi-frame aggregation \\
Unusual lighting & 2-3\% & Photometric augmentation in training \\
\bottomrule
\end{tabular}
\end{table}

\section{Code Implementation Highlights}

\subsection{LayerNorm Tuning Implementation}

\begin{lstlisting}[caption=LayerNorm-Specific Fine-tuning]
class FrameEncoder(nn.Module):
    def __init__(self, device="cpu", layernorm_tuning=True):
        # Load pretrained DINOv3
        self.model = timm.create_model(
            "vit_base_patch16_dinov3.lvd1689m",
            pretrained=True, num_classes=0
        )
        
        # Freeze all parameters
        for p in self.model.parameters():
            p.requires_grad = False
        
        # Unfreeze only LayerNorm modules
        if layernorm_tuning:
            for m in self.model.modules():
                if isinstance(m, nn.LayerNorm):
                    for p in m.parameters():
                        p.requires_grad = True
        
        # Additionally unfreeze last block
        if hasattr(self.model, 'blocks'):
            last_block = self.model.blocks[-1]
            for p in last_block.parameters():
                p.requires_grad = True
\end{lstlisting}

\subsection{Metric Loss Implementation}

\begin{lstlisting}[caption=Angular Margin Metric Loss]
def metric_loss(embeddings, labels, margin=0.6):
    """
    Compute metric loss with angular margin.
    
    Args:
        embeddings: L2-normalized embeddings [B, D]
        labels: Ground truth labels [B]
        margin: Angular margin for separation
    """
    # Compute pairwise cosine similarity
    cos_sim = F.cosine_similarity(
        embeddings.unsqueeze(1),
        embeddings.unsqueeze(0),
        dim=2
    )
    
    # Create target matrix (1 if different labels)
    targets = (labels.unsqueeze(0) != labels.unsqueeze(1)).float()
    
    # Apply margin and compute loss
    loss = (targets * F.relu(cos_sim + margin)).mean()
    
    return loss
\end{lstlisting}

\subsection{Video Prediction Pipeline}

\begin{lstlisting}[caption=Video-Level Prediction with Uncertainty]
@torch.no_grad()
def predict_video(self, frames: torch.Tensor):
    self.eval()
    frames = frames.to(self.device)
    
    # Forward pass
    logits, _ = self.forward(frames)
    probs = torch.softmax(logits, dim=-1)[:, 1]
    
    # Aggregate statistics
    per_frame_probs = probs.cpu().tolist()
    mean_score = float(probs.mean().item())
    
    # Uncertainty quantification
    instability = float(probs.std().item())
    confidence = float(abs(mean_score - 0.5) * 2.0)
    
    return {
        "score": mean_score,
        "label": "FAKE" if mean_score > 0.5 else "REAL",
        "details": {
            "per_frame_probs": per_frame_probs,
            "instability": instability,
            "confidence": confidence,
            "frame_count": len(frames)
        }
    }
\end{lstlisting}

\section{Suggested Figures for Paper}

\subsection{Figure 1: Architecture Diagram}
\textbf{Description}: Block diagram showing:
\begin{itemize}
    \item Input video $\rightarrow$ Frame extraction
    \item MTCNN face detection and cropping
    \item DINOv3 ViT encoder (with frozen blocks and unfrozen LayerNorm highlighted)
    \item Linear probe classifier
    \item Output: Real/Fake probability + uncertainty metrics
\end{itemize}

\subsection{Figure 2: Training Curves}
\textbf{Description}: Three subplots showing:
\begin{enumerate}
    \item Training loss (cross-entropy + metric loss) over epochs
    \item Validation AUROC over epochs
    \item Learning rate schedule (cosine annealing)
\end{enumerate}

\subsection{Figure 3: Embedding Visualization}
\textbf{Description}: t-SNE or UMAP projection of 768-dimensional embeddings:
\begin{itemize}
    \item Blue points: Real faces
    \item Red points: Fake faces
    \item Show clear clustering separation
\end{itemize}

\subsection{Figure 4: Comparison Bar Chart}
\textbf{Description}: Grouped bar chart comparing:
\begin{itemize}
    \item X-axis: Methods (Xception, EfficientNet, ResNet18, CLIP, DINOv3-Ours)
    \item Y-axis: AUROC scores
    \item Error bars showing variance across runs
\end{itemize}

\subsection{Figure 5: ROC Curves}
\textbf{Description}: ROC curves on Celeb-DF validation:
\begin{itemize}
    \item One curve per method
    \item AUC values annotated
    \item Diagonal reference line (random)
\end{itemize}

\subsection{Figure 6: Ablation Study Results}
\textbf{Description}: Stacked bar or waterfall chart showing:
\begin{itemize}
    \item Baseline (frozen): 0.82 AUROC
    \item + LayerNorm tuning: +0.03
    \item + Last block unfreezing: +0.02
    \item + Metric learning: +0.01+
\end{itemize}

\section{Broader Impact and Limitations}

\subsection{Ethical Considerations}

\textbf{Positive Impacts:}
\begin{itemize}
    \item Enables detection of misinformation and fraud
    \item Protects individuals from non-consensual deepfake content
    \item Supports journalistic integrity and fact-checking
\end{itemize}

\textbf{Potential Misuse:}
\begin{itemize}
    \item Could be used to filter only ``perfect'' deepfakes
    \item May create false confidence in detection capabilities
    \item Adversarial attacks could evolve to specifically target DINOv3
\end{itemize}

\subsection{Limitations}

\begin{enumerate}
    \item \textbf{Face Dependency}: Requires detectable faces; fails on videos without faces or with extreme occlusion.
    
    \item \textbf{English-Centric Pretraining}: DINOv3 trained primarily on Western faces; performance may degrade on other demographics.
    
    \item \textbf{Computational Cost}: 330MB model size may be prohibitive for edge/mobile deployment.
    
    \item \textbf{Temporal Blindness}: Frame-level analysis without explicit temporal modeling (unlike 3D CNNs or LSTMs).
    
    \item \textbf{Adaptation Lag}: May struggle with zero-day deepfake techniques not represented in training data.
\end{enumerate}

\section{Future Work}

\begin{enumerate}
    \item \textbf{Temporal Modeling}: Integrate transformer temporal attention across frames for video-level consistency checks.
    
    \item \textbf{Lightweight Deployment}: Knowledge distillation to MobileNet or EfficientNet for mobile deployment.
    
    \item \textbf{Multi-Modal Fusion}: Combine visual detection with audio analysis (lip-sync inconsistency detection).
    
    \item \textbf{Adversarial Robustness}: Train with adversarial examples to improve robustness against evasion attacks.
    
    \item \textbf{Active Learning}: Online learning framework to adapt to emerging deepfake techniques with minimal labeled data.
\end{enumerate}

\end{document}